\documentclass{article}
\usepackage[margin=1in,a4paper]{geometry}
\usepackage[utf8]{inputenc}
\usepackage[cyr]{aeguill}
\usepackage[francais]{babel}
\usepackage{hyperref}
\usepackage{amsmath}
\usepackage{gensymb}
\usepackage{enumitem,amssymb}
\newlist{checks}{itemize}{2}
\setlist[checks]{label=$\square$}
\usepackage{graphicx}
\usepackage{amsthm}
\usepackage{amsfonts}
\usepackage{multicol}
\usepackage{pgfplots}
\pgfplotsset{compat=newest}
\usetikzlibrary{calc}
\usepackage{mathtools}
\usepackage{array}
\usepackage[T1]{fontenc}
\usepackage{lmodern}
\usepackage{tabularx}
\usepackage{fancyhdr}
\usepackage{pst-func}
\usepackage{xcolor}
\usepackage{nicefrac}
\usepackage{mdframed}
\usepackage[boxed,vlined]{algorithm2e}
\usepackage{cleveref}
\newcommand{\Lim}[1]{\raisebox{0.5ex}{\scalebox{1}{$\displaystyle \lim_{#1}\;$}}}
\usepackage{float}
%\usepackage[top=2.5cm, bottom=2cm, left=2cm, right=2cm, showframe]{geometry}
\usepackage[top=2.5cm, bottom=2cm, left=2cm, right=2cm]{geometry}
\newcommand{\N}{\mathbb{N}}
\newcommand{\R}{\mathbb{R}}
\renewcommand{\C}{\mathbb{C}}
\renewcommand{\P}{\mathbb{P}}
\newcommand{\w}{\omega}
\newcommand{\p}{\partial}
\newcommand{\cross}{\times}
\newcommand{\Col}{\text{Col}}
\newcommand{\Tr}{\text{Tr}}
\newcommand{\bigzero}{\makebox(0,0){\text{\huge0}}}
\DeclareMathOperator{\Ima}{Im}
\DeclareMathOperator{\Vect}{Vect}
\usepackage{mathtools, stmaryrd}
\usepackage{xparse} \DeclarePairedDelimiterX{\Iintv}[1]{\llbracket}{\rrbracket}{\iintvargs{#1}}
\NewDocumentCommand{\iintvargs}{>{\SplitArgument{1}{,}}m}
{\iintvargsaux#1} %
\NewDocumentCommand{\iintvargsaux}{mm} {#1\mkern1.5mu..\mkern1.5mu#2}
\newcommand{\Cc}{{\mathbb{C}}}
\newcommand{\Ct}{\Cc^\times}
\newcommand{\Hh}{{\mathbb{H}}}
\newcommand{\Nn}{{\mathbb{N}}}
\newcommand{\Zz}{{\mathbb{Z}}}
\newcommand{\Zzn}{\Zz/n\Zz}
\newcommand{\ZzNt}{(\Zz/N\Zz)^\times}%quotient 
\newcommand{\Rr}{{\mathbb{R}}}
\newcommand{\Rt}{{\Rr^\times}}
\newcommand{\Qt}{{\Qq^\times}}
\newcommand{\Qq}{{\mathbb{Q}}}
% Custom titling
\usepackage{titling}
\usepackage{tcolorbox}

\title{\textbf{Analyse Avancée -- Corrigé 3}}

\renewcommand{\maketitlehookb}{
\begin{center}
\includegraphics[width=2cm]{Cours/assets/imgs/logo-round.png}
\end{center}
}

\author{Students 4 Students}
\date{Septembre 2022}

% Exercises environment and styling
\usepackage{amsthm}
\newtheoremstyle{exercice}%
    {3pt}% Space above
    {3pt}% Space below
    {\large}% Body font
    {}% Indent amount
    {\bfseries}% Theorem head font
    {}% Punctuation after theorem heading
    {\newline}% Space after theorem heading
    {\thmname{#1}\thmnumber{ #2}\thmnote{: #3}}% Theorem head spec (can be left empty, meaning ‘normal’)
\theoremstyle{exercice}
\newtheorem{exercice}{Exercice}


\begin{document}

% Header
\pagestyle{fancy}
\fancyhead[L]{Students 4 Students}
\fancyhead[C]{\textit{Analyse avancée - Corrigé 1}}
\fancyhead[R]{Septembre 2022}

\maketitle



\begin{exercice}[Convergence de séries]
Les limites des séries, si elles existent, sont les suivantes:
\begin{enumerate}
    \item On observe que $0\leq\dfrac{n}{n^3+1}\leq \dfrac{n}{n^3}=\dfrac{1}{n^2} $. En utilisant le critère de comparaison et sachant que l'on a montré dans le cours que la série de terme général $\dfrac{1}{n^2}$ converge, la série est donc convergente.
    \item Il est possible de procéder par comparaison en observant que $\dfrac{1}{n}\leq \dfrac{1}{\sqrt{n}}$. Or nous avons vu dans le cours que la série de terme général $\dfrac{1}{n}$ diverge. Ainsi, la série étudiée diverge aussi.
    \item Pour cette série, la présence d'une factorielle nous indique d'utiliser le critère du quotient. Ainsi:
    \begin{equation}
        \lim_{n\to \infty}\left| \dfrac{u_{n+1}}{u_n}\right| =\lim_{n\to \infty} \left|  \dfrac{x^{n+1}}{(n+1)!}  \dfrac{n!}{x^n}  \right |=\lim_{n\to \infty} \left|  \dfrac{x}{(n+1)} \right |=0.
    \end{equation}
    Cette série converge donc. Par ailleurs, cette série a pour limite exp($x$), la fonction exponentielle.
    
    \item Même méthode, on utilise le critère de D'Alembert. Donc: 
    \begin{equation}
        \lim_{n\to \infty}\left| \dfrac{u_{n+1}}{u_n}\right| =\lim_{n\to \infty} \left|(-1)^{n+1} \dfrac{x^{2(n+1)}}{(2(n+1))!}\dfrac{(2n)!}{(-1)^{n} x^{2n}}\right|
    \end{equation}
    \begin{equation}
        =\lim_{n\to \infty}\left| - \dfrac{x^2}{(2n+2)\cdot (2n+1} \right|=0.
    \end{equation}
    La limite étant 0 $<$ 1, la série converge. Par ailleurs, cette série converge vers $\cos (x)$. On remarque un lien vers la série de l'exponentielle qui sera vu au cours de vos cours d'analyse. 
    \item La puissance en $n$ nous indique que l'usage du critère de Cauchy est le plus censé. Ainsi:

    \begin{equation}
        \lim_{n\to \infty}\left| \left(  \left (\dfrac{n-1}{2n+1}\right )^n\right)^{\frac{1}{n}}  \right |= \lim_{n\to \infty}\left|   \left (\dfrac{n-1}{2n+1}\right )  \right |=\dfrac{1}{2}<1.
    \end{equation}
    La série converge donc bien.
\end{enumerate}
\end{exercice}


\begin{exercice}[Preuve]
    Par définition, on a  $\underset{x\in A}{\sup}f(x)$ comme étant le supremum de l'ensemble $\{f(x) ; x\in A\}$. De plus, en séparant l'ensemble $A$, il est possible de noter que $\underset{x\in A}{\sup}f(x)$ est le suprémum de $\{f(x) ; x\in B\} \cup \{f(x) ; x\in A\backslash B\}$. Ainsi, $\underset{x\in A}{\sup}f(x) = \max \{\underset{x\in B}{\sup}f(x), \underset{x\in A \backslash B}{\sup}f(x) \}$. Ainsi, il devient évident que $\underset{x\in B}{\sup}f(x) \leq \underset{x\in A}{\sup}f(x)$.

\end{exercice}

\begin{exercice}[Étude de fonctions]
Les fonctions sont définies de la façon suivante:
    \begin{enumerate}
    \item Cette fonction est définie sur tous l'ensemble des réels. De plus, elle admet comme infimum $-\infty$ et comme supremum $\infty$. Elle n'a donc pas de maximum ou de minimum.
    \item Cette fonction est définie sur tous l'ensemble des réels. De plus, elle admet comme infimum $0$ et comme supremum $\infty$. Elle n'a donc pas de maximum mais a un minimum qui est 0.
    \item Cette fonction est définie sur tous l'ensemble des réels privé de 0. De plus, elle admet comme infimum $-\infty$ et comme supremum $\infty$. Elle n'a donc pas de maximum ou de minimum.
    \item Cette fonction est définie sur tous l'ensemble des réels. De plus, elle admet comme infimum $0$ et comme supremum $\infty$. Elle n'a donc pas de maximum ou de minimum.
    \item Cette fonction est définie sur tous l'ensemble des réels privée des nombres entiers de 0 à $n$. Il s'écrit $\Rr \backslash {0, 1,...,n}$ De plus, elle admet comme infimum $-\infty$ et comme supremum $\infty$. Elle n'a donc pas de maximum ou de minimum.
\end{enumerate}
\end{exercice}


\begin{exercice}[Création de fonctions]
    Il suffit de créer une fonction définie sur un interval ouvert. En effet, considérons la fonction suivante:

\begin{equation}
    f(x)=\begin{cases}
        2x+1, x\in (0,1)\\
        2, x=0 \text{ et } x=1
    \end{cases}.
\end{equation}
Cette fonction admet comme infimum 1 et supremum 3 qui ne sont jamais atteint bien que la fonction soit bornée.

\end{exercice}


\begin{exercice}[Exercice avancée]
    
Posons la fonction $F(x)=f(x)-x$. On a que $f(a), f(b)\in [a,b] $. Ainsi, on voit que:

\begin{equation}
    \begin{cases}
        F(a)= f(a)-a \geq 0\\
        F(b)= f(b)-b\leq 0
    \end{cases}
\end{equation}
La fonction étant continue, il est possible d'utiliser le théorème des valeurs intermédiaires, vu au lycée/gymnase ou qui sera vu plus tard. Par ce théorème, on conclut qu'il existe un point $a\leq c\leq b$ tel que $F(c)=0\iff f(c)=c$.

\end{exercice}

\begin{exercice}
    
\end{exercice}


\end{document}